\typeout{NT FILE appendix_use_cases.tex}%

\chapter{Appendix B - Use Cases Specifications}
\label{ann:use_cases_appendix}

%%%% Functional Requirements
\section*{Use Cases}

%\begin{table}[H]
%    \centering
%    \renewcommand{\arraystretch}{1.5} % Adds a bit of breathing room
%    \begin{tabularx}{\textwidth}{|l|X|c|}
%        \hline
%        \textbf{ID} & \textbf{Use Case Name} & \textbf{Page} \\ \hline
%        UC-01 & Student Ranks Internship Preferences & \pageref{tab:UC-01} \\ \hline
%        UC-02 & Coordinator Validates and Provides Feedback on Proposal & \pageref{tab:UC-02} \\ \hline
%        %\rowcolor[HTML]{F9F9F9}
%        %\multicolumn{3}{|l|}{\textit{Add further use cases as needed...}} \\ \hline
%    \end{tabularx}
%    \caption{Index of Use Case Specifications}
%\end{table}

\newpage
\renewcommand{\arraystretch}{1.5} 
\begin{xltabular}{\textwidth}{|l|X|}
    \caption{Use Case Specification: Proposal Ranking} \label{tab:UC-01} \\
    
    \hline
    \textbf{Use Case ID} & UC-01 \\ \hline
    \textbf{Use Case Name} & Student Ranks Internship/Dissertation Preferences \\ \hline
    \textbf{Primary Actor} & Student \\ \hline
    \textbf{Secondary Actors} & Edition Coordinator, System \\ \hline
    \textbf{Preconditions} &
        - Student is enrolled in current edition [FR-31]. \newline
        - Proposal preference ranking period is open (not before start date, not after deadline). \newline
        - At least one proposal is available for the edition.
        \\ \hline
    \textbf{Postconditions} &
        - Student's preference ranking is saved in the system. \newline
        - Confirmation message is displayed to the student. \newline
        - Coordinator can retrieve data for seriation.
        \\ \hline
    \textbf{Main Flow} &
        1. Student logs in and navigates to the edition they are enrolled in. \newline
        2. System displays all available proposals for the edition. \newline
        3. Student views details of each proposal (objectives, requirements, company, benefits, etc.). \newline
        4. Student selects and ranks a specified number of proposals in priority order. \newline
        5. System validates the number of proposals does not exceed the limit. \newline
        6. Student submits the ranking. \newline
        7. System confirms successful submission and displays a confirmation message. \newline
        8. System stores the ranking with timestamp [NFR-AUD-1]. \\ \hline
    \textbf{Alternative Flows} &
        \textbf{A1. Post-Deadline Attempt:} \newline
        1. System prevents submission; displays "Choosing period has closed". \newline
        2. Use case terminates. \newline
        \textbf{A2. System Failure:} \newline
        1. System displays error, prompts user to try again. \newline
        2. Student's previous preference ranking remains unchanged (rollback). \\ \hline
\end{xltabular}

\newpage
\renewcommand{\arraystretch}{1.5} 
\begin{xltabular}{\textwidth}{|l|X|}
    \caption{Use Case Specification: Coordinator Validates and Provides Feedback on Proposal} \label{tab:UC-02} \\
    
    \hline
    \textbf{Use Case ID} & UC-02 \\ \hline
    \textbf{Use Case Name} & Coordinator Validates and Provides Feedback on Proposal \\ \hline
    \textbf{Primary Actor} & Edition Coordinator \\ \hline
    \textbf{Secondary Actors} & Company Representative/Professor, System, Email Service \\ \hline
    \textbf{Preconditions} &
        - A new proposal has been submitted by a Company or Professor. \newline
        - Edition Coordinator is assigned to this edition and has validation permission [FR-22]. \newline
        - The proposal review period is active.
        \\ \hline
    \textbf{Postconditions} &
        - Proposal feedback is recorded in the system. \newline
        - Submitter is notified via email and in-app notification. \newline
        - Proposal status is updated (pending clarification, accepted, or rejected). \newline
        - Audit log records the coordinator's action with timestamp.
        \\ \hline
    \textbf{Main Flow} &
        1. Edition Coordinator logs in and navigates to "Validate Proposals" page. \newline
        2. System displays list of submitted proposals awaiting validation. \newline
        3. Coordinator selects proposal to review. \newline
        4. System displays full proposal details with all submitted information. \newline
        5. Coordinator review proposal and identifies areas requiring clarification (or accepts/rejects). \newline
        6. If clarification needed: Coordinator writes specific feedback in the feedback form. \newline
        7. Coordinator selects "Request clarification" and submits [FR-22]. \newline
        8. System updates proposal status to "Awaiting Clarification". \newline
        9. System sends email to the submitter of the proposal with feedback details [FR-53]. \newline
        10. System posts in-app notification to the submitter's dashboard [FR-54]. \newline
        11. Audit log records: Coordinator name, action, feedback content, timestamp [FR-62]. \\ \hline
    \textbf{Alternative Flows} &
        \textbf{A1. Coordinator accepts proposal without changes:} \newline
        1. Step 6 is skipped. \newline
        2. Coordinator selects "Accept Proposal". \newline
        3. Proposal becomes visible to students for choosing. \newline
        4. Submitter is notified of approval via email and in-app notification. \newline
        \textbf{A2. Coordinator rejects proposal:} \newline
        1. Coordinator selects "Reject Proposal" and provides a rejection reason \newline
        2. System updates status to "Rejected". \newline
        3. Submitter is notified with rejection reason via email [FR-53]. \newline
        4. Proposal is not visible to students \newline
        \textbf{A3. Submitter fails to respond to clarification request within deadline:} \newline
        1. System sends automated reminder email 2 days before deadline. \newline
        2. After deadline, proposal automatically moves to "Deadline Exceeded" status \newline
        3. Coordinator can re-request clarification or reject. \\ \hline
\end{xltabular}

\newpage
\renewcommand{\arraystretch}{1.5} 
\begin{xltabular}{\textwidth}{|l|X|}
    \caption{Use Case Specification: System Executes Student Seriation Algorithm} \label{tab:UC-03} \\
    
    \hline
    \textbf{Use Case ID} & UC-03 \\ \hline
    \textbf{Use Case Name} & System Executes Student Seriation Algorithm \\ \hline
    \textbf{Primary Actor} & System \\ \hline
    \textbf{Secondary Actors} & Edition Coordinator, Student, Company Representative/Professor \\ \hline
    \textbf{Preconditions} &
        - The seriation period has begun (coordinator has triggered seriation) \newline
        - All students have submitted their preference rankings [FR-31] \newline
        - All accepted proposals are locked (no further modifications allowed) \newline
        - Pre-defined seriation criteria have been configured by coordinator.
        \\ \hline
    \textbf{Postconditions} &
        - Seriation algorithm completes successfully \newline
        - Assignments lists are generated for each proposal [FR-35] \newline
        - Students are notified of seriation results via email and in-app notification \newline
        - Companies/Professors receive the list of students assigned for each of their proposals \newline
        - Audit log records seriation execution details.
        \\ \hline
    \textbf{Main Flow} &
        1. Edition Coordinator logs in and navigates to "Run Seriation" \newline
        2. System displays a summary: number of students, number of proposals, seriation parameters (ECTS weight: X\%, Average Grade weight: Y\%) \newline
        3. Coordinator reviews and confirms parameters \newline
        4. Coordinator clicks "Execute Seriation" \newline
        5. System executes the seriation algorithm [FR-33]: \newline
            \begin{itemize}
                \item {For each proposal, retrieves all student preference rankings where this proposal is ranked}
                \item {Calculates a composite score for each student: (ECTSweight x studentECTS + AVGGradeweight x studentAVGGrade)}
                \item {Sorts students by composite score in descending order}
                \item {Produces a ranked list of candidates for each proposal}
            \end{itemize}
        6. System generates proposal-to-candidate assignment results \newline
        7. System sends email notifications to all students listing their seriation results [FR-34] \newline
        8. System sends email notifications to all companies and professors with their ranked candidate lists [FR-35] \newline
        9. System displays "Seriation Complete" confirmation to coordinator \newline
        10. Audit log records: seriation execution time, number of students processed, algorithm parameters used [FR-62]  \\ \hline
    \textbf{Alternative Flows} &
        \textbf{A1. Algorithm detects data inconsistency (e.g., student has invalid Average Grade)} \newline
        1. System logs error and notifies coordinator \newline
        2. Seriation halts; coordinator must resolve data issue and retry \newline
        \textbf{A2. No students ranked a particular proposal} \newline
        1. System marks that proposal with "No Candidates" status \newline
        2. Proposal is excluded from seriation results \\ \hline
\end{xltabular}

\newpage
% Consistent with UC-01 and UC-02
\renewcommand{\arraystretch}{1.5} 
\begin{xltabular}{\textwidth}{|l|X|}
    \caption{Use Case Specification: Company/Professor Responds to Coordinator Feedback and Resubmits Proposal} \label{tab:UC-04} \\
    
    \hline
    \textbf{Use Case ID} & UC-04 \\ \hline
    \textbf{Use Case Name} & Company/Professor Responds to Coordinator Feedback and Resubmits Proposal \\ \hline
    \textbf{Primary Actor} & Company Representative/Professor \\ \hline
    \textbf{Secondary Actors} & Edition Coordinator, System \\ \hline
    \textbf{Preconditions} &
        - Coordinator has requested clarification on company's/professor's proposal [UC-02, A3]. \newline
        - Company/Professor has not yet exceeded the response deadline. \newline
        - Company Representative/Professor is logged in.
        \\ \hline
    \textbf{Postconditions} &
        - Modified proposal is resubmitted. \newline
        - Coordinator is notified of resubmission. \newline
        - Proposal is returned to "Under Review" status. \newline
        - Feedback exchange is recorded in audit trail [FR-63].
        \\ \hline
    \textbf{Main Flow} &
        1. Company Representative/Professor receives email notification of requested clarification. \newline
        2. Representative/Professor logs in to the platform and navigates to "My Proposals". \newline
        3. System displays the proposal marked "Awaiting Clarification" with coordinator's feedback highlighted. \newline
        4. Representative/Professor reviews feedback and modifies proposal details (e.g., clarifies location, adds requirements). \newline
        5. Representative/Professor adds a response comment addressing each point of feedback. \newline
        6. Representative/Professor clicks "Resubmit Proposal". \newline
        7. System validates modified proposal for completeness. \newline
        8. System updates proposal status to "Under Review". \newline
        9. System sends email to Edition Coordinator with notification of resubmission [FR-54]. \newline
        10. Audit log records: resubmission timestamp, modification details, company/professor response comment [FR-63]. \\ \hline
    \textbf{Alternative Flows} &
        \textbf{A1. Company/Professor does not respond before the deadline:} \newline
        1. System sends reminder email 2 days before deadline. \newline
        2. After deadline passes, proposal is automatically marked "Deadline Exceeded". \newline
        3. Coordinator can choose to reject or grant extension. \newline
        \textbf{A2. Company's/Professor's modifications do not resolve the coordinator's concerns:} \newline
        1. Coordinator reviews resubmitted proposal (follows UC-02 main flow). \newline
        2. If still inadequate, coordinator can request additional clarification (cycle repeats). \\ \hline
\end{xltabular}


\newpage
% Consistent with previous Use Case tables
\renewcommand{\arraystretch}{1.5} 
\begin{xltabular}{\textwidth}{|l|X|}
    \caption{Use Case Specification: Edition Coordinator Creates New Edition} \label{tab:UC-05} \\
    
    \hline
    \textbf{Use Case ID} & UC-05 \\ \hline
    \textbf{Use Case Name} & Edition Coordinator Creates New Edition \\ \hline
    \textbf{Primary Actor} & Edition Coordinator \\ \hline
    \textbf{Secondary Actors} & System \\ \hline
    \textbf{Preconditions} &
        - Coordinator is authenticated and authorized as edition coordinator. \newline
        - At least one academic program and default configuration template exist.
        \\ \hline
    \textbf{Postconditions} &
        - A new edition record is created with status “Draft” or “Configured”. \newline
        - Key parameters (dates, proposal limits, seriation rules) are stored. \newline
        - Edition is available for later activation and publication.
        \\ \hline
    \textbf{Main Flow} &
        1. Coordinator accesses the “Manage Editions” area. \newline
        2. System displays a list of existing editions and an option to create a new one. \newline
        3. Coordinator selects “Create Edition”. \newline
        4. System presents a form to define edition metadata (name, academic year, degree type, description). \newline
        5. Coordinator fills in dates (proposal submission window, student ranking window, seriation execution window). \newline
        6. Coordinator defines proposal constraints (max proposals per company, per professor, per student). \newline
        7. Coordinator configures seriation parameters (e.g., average grade and ECTS weights, tie-breaking rules). \newline
        8. Coordinator saves the edition. \newline
        9. System validates the data, creates the edition in “Draft/Configured” status and confirms creation. \\ \hline
    \textbf{Alternative Flows} &
        \textbf{A1. Invalid or missing data:} \newline
        1. At step 9, validation fails (e.g., end date before start date). \newline
        2. System shows specific validation errors and highlights problematic fields. \newline
        3. Coordinator corrects the data and retries saving. \newline
        \textbf{A2. Duplicate edition:} \newline
        1. System detects an existing edition with conflicting key identifiers (e.g., same academic year and cycle). \newline
        2. System warns coordinator and suggests editing the existing edition instead. \newline
        3. Coordinator either cancels or changes identifiers and saves again. \\ \hline
\end{xltabular}

\newpage
% Consistent with previous Use Case tables
\renewcommand{\arraystretch}{1.5} 
\begin{xltabular}{\textwidth}{|l|X|}
    \caption{Use Case Specification: Company Representative Registers and Submits Initial Proposal} \label{tab:UC-06} \\
    
    \hline
    \textbf{Use Case ID} & UC-06 \\ \hline
    \textbf{Use Case Name} & Company Representative Registers and Submits Initial Proposal \\ \hline
    \textbf{Primary Actor} & Company Representative \\ \hline
    \textbf{Secondary Actors} & System, Secretariat (for protocol validation), Coordinator \\ \hline
    \textbf{Preconditions} &
        - Call for proposals for at least one edition is open. \newline
        - Representative has a valid email address.
        \\ \hline
    \textbf{Postconditions} &
        - A new company account (or user under an existing company) is created and marked as “Pending Validation” or “Active”. \newline
        - At least one proposal is stored with status “Submitted” or “Pending Validation”. \newline
        - Coordinator can see the proposal in the validation queue.
        \\ \hline
    \textbf{Main Flow} &
        1. Representative accesses the platform and chooses “Register Company / Representative”. \newline
        2. System displays a registration form (company info, contact details, legal identifiers). \newline
        3. Representative fills in the form and submits. \newline
        4. System validates data, creates the company (or links user to existing company), and sends an activation email. \newline
        5. Representative activates the account via the email link and logs in. \newline
        6. Representative navigates to “Submit Proposals”. \newline
        7. System shows available editions and their submission windows. \newline
        8. Representative selects an edition and fills in proposal data (title, description, objectives, requirements, location, compensation/benefits). \newline
        9. Representative submits the proposal. \newline
        10. System validates the proposal, stores it with status “Submitted/Pending Validation”, and shows a confirmation. \\ \hline
    \textbf{Alternative Flows} &
        \textbf{A1. Company already exists:} \newline
        1. At step 3, system detects that the company exists. \newline
        2. System offers to link the new user to the existing company instead of creating a new one. \newline
        3. Representative confirms and proceeds; account is created under that company. \newline
        \textbf{A2. Registration incomplete:} \newline
        1. Representative closes the browser or cancels before activation. \newline
        2. System keeps a temporary record, which can be completed later via emailed link or is purged after a timeout. \newline
        \textbf{A3. Proposal submitted after deadline:} \newline
        1. At step 8, if the edition’s proposal window is closed, system blocks submission. \newline
        2. System shows an error and suggests choosing another open edition (if available). \\ \hline
\end{xltabular}

\newpage
% Consistent with the rest of the Use Case Annex
\renewcommand{\arraystretch}{1.5} 
\begin{xltabular}{\textwidth}{|l|X|}
    \caption{Use Case Specification: Student Submits Support Documentation} \label{tab:UC-07} \\
    
    \hline
    \textbf{Use Case ID} & UC-07 \\ \hline
    \textbf{Use Case Name} & Student Submits Support Documentation \\ \hline
    \textbf{Primary Actor} & Student \\ \hline
    \textbf{Secondary Actors} & Coordinator, System \\ \hline
    \textbf{Preconditions} &
        - Student is authenticated and enrolled in an edition. \newline
        - Student has already created or is creating a ranking of proposals. \newline
        - Coordinator or edition configuration requires supporting documents for that edition or specific proposals.
        \\ \hline
    \textbf{Postconditions} &
        - Supporting documents (e.g., CV, motivation letter, recommendations) are uploaded and linked to the student’s candidacies. \newline
        - Coordinator and companies (where applicable) can access documents during seriation and evaluation.
        \\ \hline
    \textbf{Main Flow} &
        1. Student opens the “My Candidacies / Rankings” page. \newline
        2. System indicates which candidacies or editions require supporting documents. \newline
        3. Student selects a candidacy requiring documents. \newline
        4. System displays required and optional document types (e.g., CV required, motivation letter optional). \newline
        5. Student uploads the requested files and confirms. \newline
        6. System validates file types and sizes, stores them securely, and links them to the corresponding candidacy. \newline
        7. System shows a confirmation and updates the candidacy status to “Documents Submitted”. \\ \hline
    \textbf{Alternative Flows} &
        \textbf{A1. Missing required document:} \newline
        1. At step 5, student omits a required file. \newline
        2. System highlights the missing document and prevents completion until it is provided. \newline
        \textbf{A2. Invalid file:} \newline
        1. Uploaded file exceeds size limit or uses a forbidden type. \newline
        2. System rejects that file and informs the student of allowed types and limits. \newline
        \textbf{A3. Update documents:} \newline
        1. Before the document deadline, student returns to the candidacy. \newline
        2. System allows replacing or adding documents. \newline
        3. New versions are stored and previous versions are either archived or overwritten. \\ \hline
\end{xltabular}

\newpage
% Consistent with the rest of the Use Case Annex
\renewcommand{\arraystretch}{1.5} 
\begin{xltabular}{\textwidth}{|l|X|}
    \caption{Use Case Specification: Company/Professor Views Assigned Candidates List and Schedules Interviews} \label{tab:UC-08} \\
    
    \hline
    \textbf{Use Case ID} & UC-08 \\ \hline
    \textbf{Use Case Name} & Company/Professor Views Assigned Candidates List and Schedules Interviews \\ \hline
    \textbf{Primary Actor} & Company Representative/Professors \\ \hline
    \textbf{Secondary Actors} & Coordinator, System \\ \hline
    \textbf{Preconditions} &
        - Seriation for the relevant edition has been executed successfully. \newline
        - Company/Professor has at least one accepted proposal in that edition. \newline
        - Representative/Professor is authenticated and linked to that company.
        \\ \hline
    \textbf{Postconditions} &
        - Representative/Professor can see ranked or assigned candidates per proposal.
        \\ \hline
    \textbf{Main Flow} &
        1. Representative/Professor logs in and navigates to “My Proposals / Candidates”. \newline
        2. System lists the company’s/professor's proposals for the edition and their current seriation status. \newline
        3. Representative/Professor selects a proposal. \newline
        4. System displays the ranked candidate list or assigned students, including basic profile data allowed by policy. \newline
        5. Representative/Professor optionally selects one or more candidates get their contacts to schedule interview outside the platform. \\ \hline
    \textbf{Alternative Flows} &
        \textbf{A1. No candidates available:} \newline
        1. At step 4, system indicates that no candidates were ranked for that proposal. \newline
        2. Representative/Professor may close the view or contact the coordinator outside this use case. \\ \hline
\end{xltabular}

\newpage

\newpage
% Consistent with the rest of the Use Case Annex
\renewcommand{\arraystretch}{1.5} 
\begin{xltabular}{\textwidth}{|l|X|}
    \caption{Use Case Specification: Secretariat Verifies Payment and Finalizes Protocol} \label{tab:UC-09} \\
    
    \hline
    \textbf{Use Case ID} & UC-09 \\ \hline
    \textbf{Use Case Name} & Secretariat Verifies Payment and Finalizes Protocol \\ \hline
    \textbf{Primary Actor} & Secretariat \\ \hline
    \textbf{Secondary Actors} & Company Representative, Student, System \\ \hline
    \textbf{Preconditions} &
        - A student has been assigned to a proposal that requires a protocol or payment. \newline
        - Required protocol documents and/or payment information are available (uploaded or received externally). \newline
        - Secretariat staff is authenticated with appropriate permissions.
        \\ \hline
    \textbf{Postconditions} &
        - Payment status and protocol status are updated (e.g., “Verified”, “Signed”). \newline
        - The associated internship/dissertation is marked as “Active” or “Authorized”. \newline
        - Relevant parties can see the updated status.
        \\ \hline
    \textbf{Main Flow} &
        1. Secretariat logs in and opens the “Protocols / Payments” section. \newline
        2. System shows a list of pending cases (assigned proposals requiring verification). \newline
        3. Secretariat selects a specific case. \newline
        4. System shows student, company, proposal details, and any uploaded documents or payment records. \newline
        5. Secretariat verifies that all required documents are signed and any required payments are confirmed. \newline
        6. Secretariat updates the status to “Verified / Finalized” and saves. \newline
        7. System updates the associated assignment to “Active” and notifies the company and student. \\ \hline
    \textbf{Alternative Flows} &
        \textbf{A1. Missing or invalid documentation:} \newline
        1. At step 5, Secretariat detects missing signatures or incorrect information. \newline
        2. Secretariat sets the status to “Incomplete” and records a comment. \newline
        3. System notifies the company and/or student with the requested corrections. \newline
        \textbf{A2. Payment not received:} \newline
        1. Payment is not yet confirmed. \newline
        2. Secretariat records the issue and leaves the status as “Pending Payment”. \newline
        3. No activation occurs until payment is confirmed. \\ \hline
\end{xltabular}

\newpage
% Consistent with the rest of the Use Case Annex
\renewcommand{\arraystretch}{1.5} 
\begin{xltabular}{\textwidth}{|l|X|}
    \caption{Use Case Specification: Edition Coordinator Publishes Call for Proposals} \label{tab:UC-10} \\
    
    \hline
    \textbf{Use Case ID} & UC-10 \\ \hline
    \textbf{Use Case Name} & Edition Coordinator Publishes Call for Proposals \\ \hline
    \textbf{Primary Actor} & Edition Coordinator \\ \hline
    \textbf{Secondary Actors} & System, Company Representatives \\ \hline
    \textbf{Preconditions} &
        - At least one edition exists in “Configured” state with defined proposal submission dates. \newline
        - Coordinator is authenticated and has edition management rights. \newline
        - Company contacts exist or can be targeted by the call.
        \\ \hline
    \textbf{Postconditions} &
        - The edition’s call for proposals is marked as “Published”. \newline
        - Companies can see the call when logging in and can submit proposals for that edition. \newline
        - Notification emails are sent to targeted companies.
        \\ \hline
    \textbf{Main Flow} &
        1. Coordinator opens the “Editions” view and selects a configured edition. \newline
        2. System shows current configuration and call status. \newline
        3. Coordinator reviews details (dates, constraints) and chooses “Publish Call for Proposals”. \newline
        4. System requests confirmation and optionally allows editing the email template. \newline
        5. Coordinator confirms publication. \newline
        6. System changes the edition state to “Call Published / Open for Proposals”. \newline
        7. System sends notification emails to configured company recipients and updates the dashboard. \\ \hline
    \textbf{Alternative Flows} &
        \textbf{A1. Misconfigured edition:} \newline
        1. At step 3, system detects missing critical parameters (e.g., no submission window). \newline
        2. System blocks publication and presents a list of missing items. \newline
        3. Coordinator fixes the configuration, then retries. \newline
        \textbf{A2. Email sending issues:} \newline
        1. At step 7, some emails fail to send. \newline
        2. System logs failures and informs the coordinator, offering a way to export the list of failed addresses. \\ \hline
\end{xltabular}

\newpage
% Consistent with the rest of the Use Case Annex
\renewcommand{\arraystretch}{1.5} 
\begin{xltabular}{\textwidth}{|l|X|}
    \caption{Use Case Specification: Administrator Performs Superuser Access for Student Support} \label{tab:UC-12} \\
    
    \hline
    \textbf{Use Case ID} & UC-12 \\ \hline
    \textbf{Use Case Name} & Administrator Performs Superuser Access for Student Support \\ \hline
    \textbf{Primary Actor} & System Administrator \\ \hline
    \textbf{Secondary Actors} & Student, Coordinator, Company Representative, System \\ \hline
    \textbf{Preconditions} &
        - Administrator is authenticated with elevated permissions. \newline
        - A support request has been raised (e.g., user cannot access their account, data inconsistency, wrong role assignment). \newline
        - Superuser actions are enabled and auditable.
        \\ \hline
    \textbf{Postconditions} &
        - The specific problem (access, data correction, role fix) is resolved. \newline
        - All actions executed under superuser mode are logged with timestamp, actor, and justification.
        \\ \hline
    \textbf{Main Flow} &
        1. Administrator opens the “Support / Superuser Access” area. \newline
        2. System lists open support tickets or allows searching by user, edition, or proposal. \newline
        3. Administrator selects a case and views context (user roles, assignments, logs). \newline
        4. Based on the case the administrator logs in the Superuser with the other user's login. \newline
        5. Administrators troubleshoots the problems. \newline
        6. Logs out of the accounts being accessed. [A2] \\ \hline 
    \textbf{Alternative Flows} &
        \textbf{A1. Action not allowed by policy:} \newline
        1. The requested operation violates locked-data policy. \newline
        2. System blocks the action and explains that such changes require a different process. \newline
        3. Administrator records this outcome in the ticket. \newline
        \textbf{A2. Impersonation session ended:} \newline
        1. During impersonation, administrator finishes troubleshooting. \newline
        2. Administrator exits superuser session explicitly. \newline
        3. System returns to the administrator’s own account and logs the full session. \\ \hline
\end{xltabular}
